\ifdefined\gkver\else\def\gkver{student}\fi
\documentclass[\gkver]{gaokaozhenti}
\begin{document}

\chapter{2026年全国理综}
%% paperType: 理综物理部分

\begin{enumerate}
\item
%% number: 1
%% paperName: 2026年全国理综第 1 题 6 分
%% typeId: 1
%% type: 单选题
%% chapter: 功和能
%% point: 机械能守恒定律
%% method: 无
%% score: 6
%% degree: 750
%% duplicateId: 0
%% body:
一皮球由静止自由下落高度 $h$ 后被地面反弹，反弹后竖直向上运动。由于与地面碰撞，皮球动能损失了 $64\%$。忽略空气阻力，皮球反弹后可上升的最大高度为\xzanswer{D}

\fourchoices[answer=D]
{$0.8h$}
{$0.64h$}
{$0.6h$}
{$0.36h$}

%% answer: D
%% memo:
\memoanswer{
设落地前瞬间皮球的动能为 $E_{k}$，下落过程根据机械能守恒可得 $E_{k}=mgh$

由于与地面碰撞，皮球动能损失了 $64\%$，则皮球反弹瞬间的动能为 $E_{k}'=(1-64\%)E_{k}=36\%E_{k}$

上升过程根据机械能守恒可得 $E_{k}'=mgh'$

解得皮球反弹后可上升的最大高度为 $h'=0.36h$

故选 D。
}

\item
%% number: 2
%% paperName: 2026年全国理综第 2 题 6 分
%% typeId: 1
%% type: 单选题
%% chapter: 原子和原子核
%% point: 天然放射性
%% method: 无
%% score: 6
%% degree: 850
%% duplicateId: 0
%% body:
自由中子可以衰变，其衰变方程为 \ce{^1_0 n -> ^1_1 p + ^0_{-1} e + \overline{\nu}_e}，其中 $\overline{\nu}_{e}$ 为反中微子（质量近似为零）。由中子的衰变方程可知\xzanswer{A}

\fourchoices[answer=A]
{质子的质量小于中子的质量}
{中子的质量小于质子和电子的质量之和}
{反中微子的电荷量与质子的电荷量相等}
{反中微子的电荷量与电子的电荷量相等}

%% answer: A
%% memo:
\memoanswer{
AB．中子衰变过程释放能量，存在质量亏损，即衰变前中子的质量大于衰变后所有产物（质子、电子、反中微子，反中微子质量近似为 $0$）的质量之和，因此中子的质量大于质子与电子的质量之和，质子的质量必然小于中子的质量，故 A 正确，B 错误；

CD．根据衰变方程的电荷数守恒，左侧总电荷数为 $0$，右侧质子电荷数为 $+1$、电子电荷数为 $-1$，因此反中微子的电荷数为 $0$，既不等于质子的电荷量也不等于电子的电荷量，故 CD 错误。

故选 A。
}

\item
%% number: 3
%% paperName: 2026年全国理综第 3 题 6 分
%% typeId: 1
%% type: 单选题
%% chapter: 力物体的平衡
%% point: 三力平衡
%% method: 无
%% score: 6
%% degree: 750
%% duplicateId: 0
%% body:
吊环是男子竞技体操项目之一。双环由两吊绳分别悬挂。某运动员两手各握一环以静止方式展示了下列四种姿态。若同种姿态下两吊绳中的张力大小相等，则四种姿态下吊绳中张力最小的是\xzanswer{D}

\fourchoices[answer=D, ispicture=true, h=2.7cm]
{figs/03a.jpg}
{figs/03b.jpg}
{figs/03c.jpg}
{figs/03d.jpg}

%% answer: D
%% memo:
\memoanswer{
根据题意，设吊绳与竖直方向的夹角为 $\theta$，每根吊绳中张力为 $F$，由平衡条件有

$2F\cos\theta=mg$

解得 $F=\frac{mg}{2\cos\theta}$

可知，当 $\theta=0^{\circ}$ 时，$F$ 最小。

故选 D。
}

\item
%% number: 4
%% paperName: 2026年全国理综第 4 题 6 分
%% typeId: 1
%% type: 单选题
%% chapter: 气体性质
%% point: 气体图象
%% method: 无
%% score: 6
%% degree: 650
%% duplicateId: 0
%% body:
如图，一定量理想气体从状态 $a$ 开始，沿 $p-V$ 图上实线所示过程经状态 $b$、$c$ 到达状态 $d$。四个状态下气体内能相等的是\xzanswer{B}

\onepicture[w=5.5cm]{figs/04.png}

\fourchoices[answer=B]
{状态 $a$ 与状态 $c$}
{状态 $a$ 与状态 $d$}
{状态 $b$ 与状态 $c$}
{状态 $b$ 与状态 $d$}

%% answer: B
%% memo:
\memoanswer{
根据理想气体状态方程可得

\[ \frac{p_{a}V_{a}}{T_{a}}=\frac{p_{b}V_{b}}{T_{b}}=\frac{p_{c}V_{c}}{T_{c}}=\frac{p_{d}V_{d}}{T_{d}} \]

其中 $p_{a}=p_{b}=4\times10^{4}\UPa$，$p_{c}=p_{d}=1\times10^{4}\UPa$，$V_{a}=2\UL$，$V_{b}=V_{c}=5\UL$，$V_{d}=8\UL$

可得 $T_{b}>T_{a}=T_{d}>T_{c}$

温度相同，内能相同，可知四个状态下气体内能相等的是状态 $a$ 与状态 $d$。

故选 B。
}

\item
%% number: 5
%% paperName: 2026年全国理综第 5 题 6 分
%% typeId: 1
%% type: 单选题
%% chapter: 机械波
%% point: 波的图象
%% method: 无
%% score: 6
%% degree: 450
%% duplicateId: 0
%% body:
一列横波沿 $x$ 轴正向传播，位于坐标原点的波源振动周期为 $T$，其开始一个周期的振动图像如图所示。该波在 $t=T$ 时刻的波形图可能为\xzanswer{C}

\onepicture[w=4.2cm]{\includesvg{figs/05a}}

\fourchoices[answer=C, ispicture=true, h=2.3cm]
{figs/05b.png}
{figs/05c.png}
{figs/05d.png}
{figs/05e.png}

%% answer: C
%% memo:
\memoanswer{
从振动图像可知，$t=0$ 时波源从平衡位置向 $y$ 轴正方向起振，因此波传播到的所有质点，起振方向都沿 $y$ 轴正方向。经过一个周期 $T$，波沿 $x$ 轴正方向传播的距离为一个波长 $\lambda$，此时最前端（$x=\lambda$ 处）的质点刚要开始振动，振动方向必然沿 $y$ 轴正方向。对照一个周期的振动图像可知该波在 $t=T$ 时刻的波形图为先向下起伏、到达最低点、之后向上起伏到达最高点，且向上的起伏的最高点比向下起伏的最低点值大。

故选 C。
}

\item
%% number: 6
%% paperName: 2026年全国理综第 6 题 6 分
%% typeId: 2
%% type: 多选题
%% chapter: 运动和力
%% point: 超重失重
%% method: 无
%% score: 6
%% degree: 750
%% duplicateId: 0
%% body:
2026 年 4 月，神舟二十一号乘组的航天员在中国空间站进行了出舱活动，可认为航天员出舱后与空间站一同绕地球做圆周运动。下列说法正确的是\xzanswer{AC}

\fourchoices[answer=AC]
{出舱的航天员处于失重状态}
{出舱的航天员不处于失重状态}
{在舱内的航天员处于失重状态}
{在舱内的航天员不处于失重状态}

%% answer: AC
%% memo:
\memoanswer{
AB．出舱的航天员，绕地球做匀速圆周运动，受到地球的万有引力提供所需的向心力，加速度为重力加速度，处于完全失重状态，故 A 正确，B 错误；

CD．在舱内的航天员万有引力全部提供向心力，处于完全失重状态，故 C 正确，D 错误。

故选 AC。
}

\item
%% number: 7
%% paperName: 2026年全国理综第 7 题 6 分
%% typeId: 2
%% type: 多选题
%% chapter: 电磁感应
%% point: 验证电磁感应实验
%% method: 无
%% score: 6
%% degree: 750
%% duplicateId: 0
%% body:
利用如图所示的实验装置探究产生感应电流的条件，线圈 P 通过滑动变阻器和开关连接到电源上，线圈 Q 的两端连接到电流表上，线圈 P 装在线圈 Q 的里面。下列说法正确的是\xzanswer{AC}

\onepicture[w=6.5cm]{figs/07.png}

\fourchoices[answer=AC]
{开关闭合瞬间，线圈 Q 中产生感应电流}
{开关断开瞬间，线圈 Q 中不产生感应电流}
{开关闭合状态下，快速滑动变阻器滑片，线圈 Q 中产生感应电流}
{开关闭合状态下，快速滑动变阻器滑片，线圈 Q 中不产生感应电流}

%% answer: AC
%% memo:
\memoanswer{
A．开关闭合瞬间，线圈 P 的电流从无到有，产生的磁场也从无到有，穿过线圈 Q 的磁通量突然增加，满足感应电流产生条件，Q 中会产生感应电流，A 正确；

B．开关断开瞬间，线圈 P 的电流从有到无，产生的磁场也从有到无，穿过线圈 Q 的磁通量突然减小，Q 中会产生感应电流，B 错误；

CD．开关闭合状态下，快速滑动变阻器滑片时，滑动变阻器接入电路的阻值快速变化，线圈 P 的电流随之快速变化，它产生的磁场也发生变化，穿过 Q 的磁通量随之改变，因此线圈 Q 中会产生感应电流，C 正确，D 错误。

故选 AC。
}

\item
%% number: 8
%% paperName: 2026年全国理综第 8 题 6 分
%% typeId: 2
%% type: 多选题
%% chapter: 电场
%% point: 带电粒子的偏转
%% method: 无
%% score: 6
%% degree: 750
%% duplicateId: 0
%% body:
两带电平行金属板竖直放置，一带电小球从高于金属板的位置由静止自由下落，以一定的速度进入两板间。若忽略边缘效应，下列表示小球运动轨迹（图中实线）的四幅图中，可能正确的是\xzanswer{AB}

\fourchoices[answer=AB, ispicture=true, h=2.7cm]
{figs/08a.png}
{figs/08b.png}
{figs/08c.png}
{figs/08d.png}

%% answer: AB
%% memo:
\memoanswer{
小球进入竖直平行板电场前做自由下落运动，进入电场瞬间的初速度方向是竖直向下的。进入电场后，小球同时受两个恒力：竖直向下的重力、水平方向的恒定电场力，因此合外力是恒定的恒力。由于进入电场瞬间水平分速度为 $0$，轨迹在电场边界处的切线方向应与进入前的速度方向一致（即竖直向下），轨迹应是光滑的。竖直方向：初速度向下，同时重力提供向下的加速度，竖直分速度会持续增大；水平方向：初速度为 $0$，在恒定电场力作用下做初速度为 $0$ 的匀加速直线运动，水平分速度从 $0$ 开始逐渐增大。因此小球的运动轨迹是曲线且向电场力方向偏转，由于不知道两金属板的电性，所以 AB 符合题意，CD 不符合题意。

故选 AB。
}

\item
%% number: 9
%% paperName: 2026年全国理综第 9 题 8 分
%% typeId: 4
%% type: 实验题
%% chapter: 曲线运动
%% point: 研究平抛运动实验
%% method: 无
%% score: 8
%% degree: 750
%% duplicateId: 0
%% body:
某同学利用如图所示的装置探究平抛运动竖直分运动的规律。

\begin{enumerate}
	\item
	他将装置固定在距离水平地面一定高度处，将球 A 夹在弹性金属片 P 和左侧固定架之间，与球 A 相同的球 B 放在固定于 P 底端的小平台上，并使两球高度相同。用小锤击打 P，球 B 水平抛出，同时球 A 自由下落，比较两球落地时间。
	\item
	保持其他条件不变，适当增加击打力度，球 B 被抛出时的速度 \tkanswer[2cm]{增大}（填“增大”“减小”或“不变”）。
	\item
	分别改变小锤击打力度和装置距地面的高度，都能得到两球同时落地的结果，由此得到了平抛运动竖直分运动与自由落体运动 \tkanswer[2.5cm]{一致}（填“一致”或“不一致”）的结论。
	\item
	假设某次实验发现球 A 比球 B 早落地，这可能是此次实验中球 B 抛出方向 \tkanswer[2.5cm]{斜向上}（填“斜向上”或“斜向下”）导致的。
\end{enumerate}

\onepicture[h=4.5cm, align=right]{\includesvg{figs/09}}

%% answer:
\jdanswer{
\begin{enumerate}
	\item
	两球同时落地。

	\item
	增大

	\item
	一致

	\item
	斜向上
\end{enumerate}
}

%% memo:
\memoanswer{
（2）弹性金属片的弹性形变程度由击打力度决定，击打力度越大，金属片获得的弹性势能越大，转化给球 B 的动能越多，因此球 B 被抛出时的水平初速度增大。

（3）改变抛出初速度、改变下落高度后两球始终同时落地，说明平抛运动在竖直方向的运动规律，和相同高度下自由落体的运动规律完全相同，因此二者竖直分运动一致。

（4）若球 B 抛出方向斜向下，它的竖直分初速度向下，竖直方向运动比自由下落的球 A 更快，会比 A 先落地；若球 B 斜向上抛出，它的竖直分初速度向上，会先向上运动再下落，总运动时间会比自由下落的 A 更长，因此会出现球 A 比球 B 早落地的情况。
}

\item
%% number: 10
%% paperName: 2026年全国理综第 10 题 18 分
%% typeId: 4
%% type: 实验题
%% chapter: 电路
%% point: 电路实验
%% method: 无
%% score: 18
%% degree: 550
%% duplicateId: 0
%% body:
某同学测量两个电压表 $V_{1}$ 和 $V_{2}$ 的内阻。可用器材有：直流电源 $E$（电源电动势未知），滑动变阻器 $R_{0}$，电阻箱 $R$，开关 S 和导线若干。

\threepicture[numA=a, numB=b, numC=c, labelA=2026全国理综10a, labelB=2026全国理综10b, labelC=2026全国理综10c, h=4.0cm]
{figs/10a.png}{figs/10b.png}{figs/10c.png}

\begin{enumerate}
	\item
	他设计了测量电压表 $V_{2}$ 内阻的电路，其电路原理图如图（a）所示。在答题卡上按照原理图将图（b）中实物图间的连线补充完整\tkanswer[2cm]{}。闭合开关 S 前，要将滑动变阻器的滑片滑至与图（a）中滑动变阻器的 \tkanswer[1.8cm]{$M$}（填“$M$”或“$N$”）端对应的位置。
	\item
	闭合开关 S，将滑动变阻器滑片移动到适当位置。增加电阻箱接入电路的阻值，$V_{1}$ 的示数发生了变化，其示数 \tkanswer[1.8cm]{增大}（填“增大”或“减小”），同时 $V_{2}$ 的示数 \tkanswer[1.8cm]{减小}（填“增大”“减小”或“不变”）。
	\item
	若电阻箱接入电路的阻值为 $R$ 时，$V_{1}$ 和 $V_{2}$ 的示数分别为 $U_{1}$ 和 $U_{2}$，则电压表 $V_{2}$ 的内阻值 $R_{V}=$ \tkanswer[3cm]{$\frac{U_2}{U_1-U_2}R$}（用 $U_{1}$、$U_{2}$ 和 $R$ 表示）。当 $R=1.4000\UkO$ 时，$V_{1}$ 的示数 $U_{1}=2.75\UV$，$V_{2}$ 的示数如图（c）所示，可知 $U_{2}=$ \tkanswer[1.8cm]{$1.35$} $\UV$，$R_{V}=$ \tkanswer[1.8cm]{$1.35$} $\UkO$（保留 2 位小数）。
	\item
	若使用前 $V_{2}$ 已调零，而 $V_{1}$ 未调零，导致 $U_{1}$ 偏大，则（3）中计算得到的电压表 $V_{2}$ 的内阻值 \tkanswer[1.8cm]{小于}（填“大于”“小于”或“等于”）$V_{1}$ 调零后相应的计算结果。
	\item
	将上述实验中两电压表的位置互换，测量得到电压表 $V_{1}$ 的内阻。
\end{enumerate}

%% answer:
\jdanswer{
\begin{enumerate}
	\item
	实物图连接见详解，$M$

	\item
	增大，减小

	\item
	$\frac{U_2}{U_1-U_2}R$，$1.35$，$1.35$

	\item
	小于

	\item
	将 $V_{1}$ 与 $V_{2}$ 位置互换后，按上述同样的方法操作，即可测量得到电压表 $V_{1}$ 的内阻。
\end{enumerate}
}

%% memo:
\memoanswer{
（1）根据电路图连接实物图

\onepicture[w=6.0cm]{figs/10_an.png}

为保护电路，且滑动变阻器采用分压式接法，满足测量电路部分电压从零开始条件，则闭合开关 S 前，滑动变阻器的滑片应置于 $M$ 端。

（2）闭合开关 S，将滑动变阻器滑片移动到适当位置。增加电阻箱接入电路的阻值，电路的总电阻增大，总电流 $I_{\text{总}}$ 减小，根据闭合电路欧姆定律并联电路两端 $U_{1}=E-I_{\text{总}}(r+R_{\text{滑右}})$ 增大，所以 $V_{1}$ 的示数增大

对滑动变阻器左侧电阻，$U_{1}=I'R_{\text{滑左}}$，可得 $I'$ 增大

根据电流关系 $I_{\text{总}}=I'+I$，$I'$ 增大，$I_{\text{总}}$ 减小，所以 $I$ 减小

\onepicture[w=4.5cm]{figs/10_an2.jpg}

对电压表 $V_{1}$ 支路，$I_{1}=\frac{U_{1}}{R_{V2}}$，$R_{V2}$ 不变，$U_{1}$ 增大，所以 $I_{1}$ 增大

根据电流关系 $I=I_{1}+I_{2}$，$I$ 减小，$I_{1}$ 增大，所以 $I_{2}$ 减小

对于电压表 $V_{2}$ 所在支路，$U_{2}=I_{2}R_{V2}$，因 $R_{V2}$ 不变，$I_{2}$ 减小，所以 $U_{2}$ 减小，即 $V_{2}$ 的示数减小

（3）$V_{2}$ 和电阻箱 $R$ 串联，电流相等，串联电路电压与电阻成正比 $\frac{U_{2}}{R_{V}}=\frac{U_{1}-U_{2}}{R}$

整理得电压表 $V_{2}$ 的内阻值为 $R_{V}=\frac{U_{2}}{U_{1}-U_{2}}R$

根据图（c）可知电压表 $V_{2}$ 的示数为 $U_{2}=1.35\UV$

根据 $R_{V}=\frac{U_{2}}{U_{1}-U_{2}}R$ 可得 $R_{V}=\frac{1.35}{2.75-1.35}\times1.4000\UkO=1.35\UkO$

（4）$V_{1}$ 未调零导致 $U_{1}$ 测量值偏大，$U_{2}$ 是准确值，因此 $U_{1}-U_{2}$ 的计算值比真实值偏大，代入 $R_{V}=\frac{U_{2}}{U_{1}-U_{2}}R$ 可得 $R_{V}$ 结果小于调零后的真实计算结果。
}

\item
%% number: 11
%% paperName: 2026年全国理综第 11 题 16 分
%% typeId: 6
%% type: 计算题
%% chapter: 曲线运动
%% point: 平抛运动
%% method: 无
%% score: 16
%% degree: 650
%% duplicateId: 0
%% body:
某同学从湖边山顶以 $v_{0}=15\Ums$ 的速度向湖中水平抛出一石块，经过 $4\Us$ 石块落到水面，再经过 $0.3\Us$ 石块撞击水面的声音传回到抛出点。若不计空气阻力，重力加速度大小取 $g=10\Umsq$，求

\begin{enumerate}
	\item
	石块抛出点距湖面的高度；
	\item
	石块落水点与抛出点之间的距离；
	\item
	声速的大小。
\end{enumerate}

%% answer:
\jdanswer{
\begin{enumerate}
	\item
	$h=80\Um$

	\item
	$s=100\Um$

	\item
	$v=\frac{1000}{3}\Ums$
\end{enumerate}
}

%% memo:
\memoanswer{
（1）设石块从抛出到落水所需时间为 $t$，抛出点距湖面的高度为 $h$，由运动学公式有 $h=\frac{1}{2}gt^{2}$

代入题给数据得 $h=80\Um$

（2）设石块落水点到抛出点的距离为 $s$，石块的水平位移为 $l$，根据运动学公式和几何关系有 $l=v_{0}t$

$s=\sqrt{l^{2}+h^{2}}$

联立数据可得 $s=100\Um$

（3）设声速的大小为 $v$，石块落水声传到石块抛出点所需的时间为 $t'$，则 $s=vt'$

联立解得 $v=\frac{1000}{3}\Ums$
}

\item
%% number: 12
%% paperName: 2026年全国理综第 12 题 20 分
%% typeId: 6
%% type: 计算题
%% chapter: 电磁感应
%% point: 电磁感应中的匀变速
%% method: 无
%% score: 20
%% degree: 350
%% duplicateId: 0
%% body:
光滑水平桌面（纸面）上，两平行虚线之间有方向垂直纸面向里的匀强磁场，一正方形导体线框 $MNPQ$ 静置在桌面上，$PQ$ 边与磁场边缘重合，如图（a）所示。$t=0$ 时刻，线框在拉力作用下由静止开始进入磁场，拉力方向平行于桌面且与虚线垂直，线框完全进入磁场后，继续运动一段时间，直至 $MN$ 边离开磁场为止。拉力大小随线框速度大小变化的 $F-v$ 图线如图（b）中实线所示，线框加速度始终保持不变，$PQ$ 边始终与磁场边缘保持平行。线框质量为 $m$、电阻为 $R$，图（b）中 $F_{0}$ 与 $v_{0}$ 已知，求

\twopicture[numA=a, numB=b, labelA=2026全国理综12a, labelB=2026全国理综12b, h=4.5cm]
{figs/12a.png}{figs/12b.png}

\begin{enumerate}
	\item
	线框加速度大小及线框边长；
	\item
	磁感应强度大小及磁场区域的宽度；
	\item
	从 $PQ$ 边进入磁场到 $MN$ 边离开磁场的过程中，拉力的最大值及安培力冲量的大小。
\end{enumerate}

%% answer:
\jdanswer{
\begin{enumerate}
	\item
	$\frac{F_0}{m}$，$\frac{mv_0^2}{2F_0}$

	\item
	$\frac{2F_0}{mv_0^2}\sqrt{\frac{2F_0R}{v_0}}$，$\frac{25mv_0^2}{8F_0}$

	\item
	$(\sqrt{29}+1)F_{0}$，$2mv_{0}$
\end{enumerate}
}

%% memo:
\memoanswer{
（1）设线框加速度大小为 $a$，由图（b）可知当线框速度大小为 $v_{0}$ 时线框完全进入磁场，整个线框在磁场中运动时所受安培力为零，合力即拉力，大小为 $F_{0}$。根据牛顿第二定律得 $F_{0}=ma$

可得线框加速度大小 $a=\frac{F_{0}}{m}$

设线框边长为 $l$，根据运动学公式有 $v_{0}^{2}=2al$

联立解得线框边长 $l=\frac{mv_{0}^{2}}{2F_{0}}$

（2）图（b）斜线部分对应的运动过程中，线框所受安培力不为零，设相应过程中线框速度的大小为 $v$ 时，感应电动势为 $E$，感应电流为 $I$，安培力大小为 $f$。设磁感应强度大小为 $B$，则产生的电动势为 $E=Blv$

由闭合电路欧姆定律 $I=\frac{E}{R}$

线框做匀加速运动，合外力大小为 $F_{0}$，拉力大小为 $F$，则 $F_{0}=F-f$

其中 $f=BIl$

联立相关各式得 $F=F_{0}+\frac{B^{2}m^{2}v_{0}^{4}v}{4F_{0}^{2}R}$

由图（b）知，$F=3F_{0}$ 时 $v=v_{0}$，代入上式解得磁感应强度大小为 $B=\frac{2F_{0}}{mv_{0}^{2}}\sqrt{\frac{2F_{0}R}{v_{0}}}$

设 $PQ$ 边到达磁场上边缘时线框速度大小为 $v_{1}$，由图（b）知，$F=6F_{0}$ 时 $v=v_{1}$，可得

$6F_{0}=F_{0}+\frac{B^{2}m^{2}v_{0}^{4}v_{1}}{4F_{0}^{2}R}$

联立可得 $v_{1}=\frac{5}{2}v_{0}$

设磁场区域宽度为 $L$，则 $v_{1}^{2}=2aL$

联立相关各式得磁场区域的宽度为 $L=\frac{25mv_{0}^{2}}{8F_{0}}$

（3）设 $MN$ 边到达磁场上边缘时线框速度大小为 $v_{2}$，此时拉力最大，根据运动学公式得 $v_{2}^{2}-v_{1}^{2}=2al$，其中 $v_{1}=\frac{5}{2}v_{0}$

联立相关各式得 $v_{2}=\frac{\sqrt{29}}{2}v_{0}$

将 $v=v_{2}$ 时 $F=F_{\max}$ 代入

$F=F_{0}+\frac{B^{2}m^{2}v_{0}^{4}v}{4F_{0}^{2}R}$

其中 $B=\frac{2F_{0}}{mv_{0}^{2}}\sqrt{\frac{2F_{0}R}{v_{0}}}$

联立可得 $F$ 的最大值 $F_{\max}$ 为 $F_{\max}=(\sqrt{29}+1)F_{0}$

设整个过程安培力的冲量大小为 $I_{f}$，由于线框始终做匀加速直线运动，运动时间 $t=\frac{v}{a}$

则 $F$ 的冲量大小等于图（b）中实线下的面积 $S$ 除以 $a$，根据动量定理得 $mv_{2}=\frac{S}{a}-I_{f}$

\onepicture[w=5.5cm]{figs/12_an.jpg}

由上图得出面积 $S=S_{1}+S_{2}+S_{3}$

再联立相关各式得 $I_{f}=2mv_{0}$
}

\end{enumerate}

\end{document}
